\section{Related work}
\label{sec:related}

\subsection{Temporal shifting}

Deferring flexible work to hours of lower carbon intensity was examined
systematically by Wiesner et al.~\cite{wiesner2021letswait}, who characterised
delay-tolerant workloads and quantified the potential for temporal shifting in
four regions over a year. Their finding that the achievable reduction depends
strongly on the region examined is the starting point for
Section~\ref{sec:res-home}: we ask which property of the region is responsible,
and find that it is the region's mean intensity rather than its temporal
structure. Earlier work on scheduling around on-site renewable supply
\cite{goiri2011greenslot,grange2018greenit} addresses a related problem, though
with generation local to the datacentre rather than grid intensity as the
varying quantity.

\subsection{Spatial and spatio-temporal shifting}

Relocating work between grids has been studied both as load balancing across
geo-distributed datacentres \cite{zhou2013carbonaware} and, more recently, as
carbon-aware placement in production settings
\cite{radovanovic2023carbonaware,souza2023casper}. Sukprasert et
al.~\cite{sukprasert2023spatiotemporal,sukprasert2024limitations} examine the
limits of both temporal and spatial shifting and conclude that reported benefits are sensitive to the
assumptions under which they are measured. Our work is complementary: where they
establish that the assumptions matter in aggregate, we isolate three specific
parameters and quantify each separately, holding the remainder fixed.

Zanotto et al.~\cite{zanotto2025usercentric} evaluate combined spatial and
temporal shifting for non-preemptible virtual machines under a Kubernetes-based
architecture. They formulate placement as a binary integer program solved one
request at a time against a greedily updated occupancy matrix, and evaluate it in
CloudSim Plus against a carbon-agnostic round-robin baseline, with a per-machine
power profile drawn from SPECpower and virtual machine requests sampled from a
production trace. Their formulation is the direct antecedent of the policies
compared in Section~\ref{sec:results}, and we adopt their host platform,
simulator, baseline, per-region concurrency constraint and deadline-margin levels
so that our results and theirs differ in as few respects as possible.

Their constraint limiting each region to $M_j$ simultaneous jobs is introduced
explicitly to prevent what they term the ``black hole'' case, in which the
optimiser assigns the entire workload to the single cleanest region. They
observe that the limit governs the outcome, that a permissive limit concentrates
nearly all work in one region while a restrictive one distributes it, and that
the value of $M_j$ must be determined experimentally. They leave that estimation
as future work. Sections~\ref{sec:res-capacity} and~\ref{sec:res-regionset}
address it, and find that the parameter accounts for a larger share of the
variation in reported savings than the scheduling policy does.

\subsection{Carbon intensity data and its representation}

The choice of carbon intensity signal is itself consequential. Maji et
al.~\cite{maji2024greenmirage} show that location-based and market-based
estimation methods can lead to different conclusions about the efficacy of the
same optimisation. Forecast-based scheduling depends further on prediction
quality, for which day-ahead \cite{maji2022dacf} and multi-day
\cite{maji2022carboncast} methods have been developed. We use published
lifecycle emission factors throughout and plan against ground truth, so that
the parameters under study are not confounded with prediction error; the effect
of forecast quality is a separate question that this paper deliberately does not
address. We note that Zanotto et al.~\cite{zanotto2025usercentric} report
closely comparable results under forecast and under perfect foresight,
which supports treating prediction error as separable from the parameters
studied here.

Infrastructure for carbon-aware experimentation has also received attention.
Testbeds such as Vessim \cite{wiesner2023vessim} and virtualised energy systems
such as Ecovisor \cite{souza2023ecovisor} provide environments in which
carbon-aware applications can be developed and evaluated. Our aim differs: we do
not propose a new environment but examine what an evaluation conducted in such
an environment is measuring.

\subsection{Simulation and power modelling}

CloudSim \cite{calheiros2011cloudsim} and its successor CloudSim Plus
\cite{silvafilho2017cloudsimplus} are widely used for cloud resource-management
studies, and we adopt the latter. Reported energy figures in such studies depend
on the host power model, and Ismail and Materwala \cite{ismail2020serverpower}
survey the models in use and their accuracy. We use a measured
SPECpower\_ssj2008 curve for the specific host modelled
\cite{spec2024xr8620t} and, because the choice is consequential, repeat the
entire evaluation under a linear approximation to establish which of our
conclusions depend on it (Section~\ref{sec:res-robustness}).

\subsection{Position of this work}

Prior work has proposed carbon-aware scheduling policies and has examined their
limits. What has not been established is how much of the variation in reported
savings is attributable to the evaluation configuration rather than to the
policies themselves. This paper addresses that question directly, and the answer
bears on how such evaluations should be reported.
